\section{TitleBlock.sty and Hyperlinks.sty}\label{sec:ref-titleblock}

Two small base packages, both built on the same hook pattern: plain
defaults here, a class overrides with \cs{renewcommand}.

\subsection*{TitleBlock}

The shared title block for LectureNotes and CourseDocument -- title,
subtitle, ``Updated $\langle$date$\rangle$'' -- structure written once,
class differences reduced to hooks. (Assessment deliberately does not load
it: an assessment opens with \cs{AssessmentTitle}, a different artefact.)
Also home of \file{datetime2}, the date's formatter, loaded
\texttt{[en-GB,calc]}.

Both of those matter beyond \cs{@date}. \texttt{ord=raise} is set here, so
\emph{every} date the toolkit sets gets a superscripted ordinal suffix ---
\emph{3\textsuperscript{rd}}, not \emph{3rd} --- without a document ever
saying \cs{textsuperscript}. And \texttt{calc} loads the module that
computes a day-of-week index, which is what lets a document name a weekday:

\begin{manexsource}
\DTMsavedate{withdrawal}{2026-11-03}
...the last day to withdraw is
  {\DTMsetup{showdow}\DTMusedate{withdrawal}}.   % Tuesday 3rd November, 2026
\end{manexsource}

The braces are load-bearing: \texttt{showdow} is a boolean, so grouping it
keeps the weekday on that date and off the \emph{Updated} line. The option
has to be passed \emph{here} rather than by the document, because datetime2
is loaded during \cs{documentclass} and disables its own \texttt{calc} key
straight afterwards --- there is no way to ask for it late. It is opt-in
upstream because it costs \file{pgfkeys} $+$ \file{pgfcalendar}; here it
costs \file{pgfcalendar} alone, since both classes loading this file already
carry \file{pgfkeys} via \file{tcolorbox}. Assessment pays nothing, loading
neither this file nor datetime2.

\begin{center}
\begin{tabularx}{\linewidth}{@{}l Y@{}}
\cs{subtitle}\texttt{\{...\}} & the subtitle line(s); set before
  \cs{maketitle}.\\
\cs{TitleColor} & a \emph{style} hook (e.g.\
  \cs{color}\texttt{\{HeadingColor\}}), deliberately not a colour name --
  it keeps TitleBlock free of any xcolor dependency. Default: empty.\\
\cs{TitleStamp} & follows the date (CourseDocument sets ``at
  $\langle$time$\rangle$''). Default: empty.\\
\cs{TitleMarker} & follows the stamp; currently a free hook -- the
  instructor-copy signal moved to the page frame. Default: empty.\\
\end{tabularx}
\end{center}

\subsection*{Hyperlinks}

hyperref, plus \cs{demo}. The load itself is the point:
hyperref patches much of LaTeX, so anything loaded after it can silently
undo those patches -- Hyperlinks therefore defers it with
\cs{AtEndPreamble}, which runs after every package the \emph{document}
loads too, not merely after the class
(section~\ref{sec:load-order}). Requires the class to have passed
\texttt{dvipsnames} to xcolor beforehand.

\begin{center}
\begin{tabularx}{\linewidth}{@{}l Y@{}}
\cs{LinkColor} & internal links (default \texttt{black}; CourseDocument
  sets its identity colour). Read late, so overrides anywhere in the
  preamble win.\\
\cs{UrlColor} & external URLs (default \texttt{RoyalBlue}).\\
\cs{secref}\texttt{\{label\}} & a cross-reference by \emph{name}, not
  number: ``the Exemptions section'', linked on the title. Numbers belong
  on headings and in the table of contents, where they index; in a
  sentence they make the reader hold a coordinate and go hunting.\\
\cs{SecRefFormat}\texttt{\{linked title\}} & that wording, in one hook
  (default \texttt{the \#1 section}). Note ``the'' lives \emph{inside} it,
  so call sites write \cs{secref} bare --- ``see the \cs{secref}\ldots''
  renders ``the the''.\\
\cs{demo}\texttt{[tool]\{URL\}} & a link out to an interactive tool; see
  below.\\
\cs{DemoLabel}, \cs{DemoMarker} & the student-facing wording
  (\texttt{interactive graph}) and its cue ($\triangleright$).\\
\cs{DemoFormat}\texttt{\{URL\}} & the whole rendering, in one hook.\\
\cs{interactive}\texttt{[dir]\{slug\}} & a \cs{demo} with a catalogue
  entry behind it: thumbnail, label and caption; see below.\\
\cs{InteractiveMeta}\texttt{\{keys\}} & a catalogue entry's metadata;
  store-only outside \cs{interactive}.\\
\cs{InteractiveDir} & the default \texttt{dir}
  (\texttt{../interactives/}, right for a lecture).\\
\cs{InteractiveThumbWidth}, \cs{InteractiveGap} & thumbnail width
  (default \texttt{0.32\textbackslash linewidth}) and the gap after it.\\
\cs{InteractivePendingNote} & the instructor-copy note under a pending
  entry.\\
\cs{InteractiveFormat} & the whole rendering, in one hook.\\
\end{tabularx}
\end{center}

\subsection*{\texorpdfstring{\cs{demo}}{\textbackslash demo}}

\emph{Students do not care which tool it is; only the author does.} A
class that meets ``LINK: DESMOS'', ``Desmos demo'' and ``LET'S EXPLORE
WITH DESMOS!'' on three consecutive pages is reading three different
things when it is always the same thing: press here, a graph moves. So
every demo prints \emph{one} uniform label:

\begin{manexsource}
\demo{https://www.desmos.com/calculator/xxxxxxxxxx}
\demo[geogebra]{https://example.org/some/opaque/url}
\end{manexsource}

The optional argument is \emph{accepted and discarded} -- the same
store-only idiom as ProblemMeta -- naming the tool for a human reading
the source, worth giving when the URL's own domain does not already say.
Nothing about it reaches the page.

\textbf{Write the whole URL, scheme and all.} A URL without
\texttt{://} is a \emph{relative} one: the viewer resolves it against the
document's own location -- a local file path -- and the link silently does
nothing, while the page still looks perfectly right. Pasting
\texttt{www.example.com} without its \texttt{https://} is the entire failure,
so \cs{demo} checks and warns at build time rather than letting it through:

\begin{manexsource}
Package Hyperlinks Warning: \demo URL has no scheme and will
(Hyperlinks)                not be clickable: www.cbc.ca on input line 29.
\end{manexsource}

Same principle as TextRef's unresolved keys (section~\ref{sec:ref-textref}):
degrade if you must, but never silently.

The real payoff is enumerability: with every interactive link routed
through one macro, ``curl all my demos and tell me which died over the
summer'' is a grep away. That is reason enough to write \cs{demo} rather
than the bare \cs{href} you were about to.

It lives in Hyperlinks rather than TextRef (section~\ref{sec:ref-textref})
deliberately: a demo has nothing to do with the textbook -- it is as often
your own Desmos page -- and Hyperlinks is loaded by LectureNotes and
CourseDocument alike, so a syllabus gets \cs{demo} too. Assessment loads
no hyperref and therefore has none.

\subsection*{\texorpdfstring{\cs{interactive}}{\textbackslash interactive}}

A bare URL cannot be searched before you build a graph that already
exists, cannot say where it came from, and cannot tell a student what
they are about to see. \cs{interactive} adds the three, from a
\emph{catalogue entry}: one directory per graph, in the topic repo beside
\file{figures/} and \file{problems/}.

\begin{manexsource}
<topic>/interactives/evt-closed-interval/meta.tex       the metadata
<topic>/interactives/evt-closed-interval/thumbnail.png  exported from the tool
\end{manexsource}

\file{meta.tex} holds one block, store-only like \cs{ProblemMeta} -- it
compiles to nothing anywhere but inside \cs{interactive}, which reads
\texttt{url} and \texttt{shows} from it. Brace any value with a comma in it;
a misspelt key stops the build.

\begin{manexsource}
\InteractiveMeta{
  url     = {https://www.desmos.com/calculator/xxxxxxxxxx},
  tool    = {desmos},
  source  = {own},
  shows   = {Drag the endpoints of the interval: on a closed interval
             the maximum and minimum are always there.},
  tags    = {extreme-values},
  section = {3.1},
}
\end{manexsource}

In a lecture, \cs{interactive}\texttt{\{evt-closed-interval\}} prints the
thumbnail on the left (linked), and on the right the same
\emph{interactive graph} label \cs{demo} prints, with the \texttt{shows}
line under it as a caption. The uniform label still keeps the tool's
identity off the page; the caption is new, and says \emph{what} moves.
The optional argument is the catalogue directory, written literally and
resolved against the calling file like \cs{importproblem}'s; the default,
\texttt{../interactives/}, suits a lecture one level below its topic root.

\textbf{A pending entry builds cleanly.} Entries are usually written before
their graph exists, and a lecture citing one must still build with zero
warnings. So an empty \texttt{url} or a missing \file{thumbnail.png} is not
a warning: a plain frame stands in for the thumbnail, the label prints
unlinked, and the instructor copy alone says \emph{pending}. A slug with
no \file{meta.tex} at all \emph{is} a warning -- that is a typo, never a
plan -- and a URL without a scheme gets \cs{demo}'s warning.
